%%%%%%%%%%%%%%%%%%%%%%%%%%%%%%%%%%%%%%%%%
% Medium Length Professional CV
% LaTeX Template
%
% Modified for GAFA-oriented resume
%
%%%%%%%%%%%%%%%%%%%%%%%%%%%%%%%%%%%%%%%%%

\documentclass[
    11pt, % Font size can be 10pt, 11pt, or 12pt
]{resume} % Use the resume class

\usepackage{ebgaramond} % Use the EB Garamond font
\usepackage[dvipdfm,linkcolor=blue,anchorcolor=blue,urlcolor=blue,colorlinks=true]{hyperref}
% \usepackage[dvipdfmx]{hyperref}
\usepackage{pxjahyper}

%----------------------------------------------------------------------------------------
%   NAME AND CONTACT INFORMATION
%----------------------------------------------------------------------------------------

\name{Kan Kusakabe} % Name to appear at the top

\address{Hokkaido University, Japan}
\address{kusakabe.kan.v5@elms.hokudai.ac.jp}
\address{\url{https://kankusakabe.github.io/Portfolio}}

\begin{document}

%----------------------------------------------------------------------------------------
%   SUMMARY SECTION (OPTIONAL)
%----------------------------------------------------------------------------------------
\begin{rSection}{Profile / Objective}
Ph.D. Candidate in Computer Science with a strong background in Human-Computer Interaction (HCI), software development, and data analysis. Skilled in rapid prototyping, UX research, machine learning, sensing, and video creation. Seeking opportunities in \textbf{Software Engineering}, \textbf{Data Science}, or \textbf{UX Research} roles, leveraging research-driven problem solving and practical development expertise. (\href{https://github.com/KanKusakabe/}{Github})
\end{rSection}

%----------------------------------------------------------------------------------------
%   TECHNICAL STRENGTHS SECTION
%----------------------------------------------------------------------------------------

\begin{rSection}{Technical Strengths}
\textbf{Domain}:
Rapid Prototyping, 3D Printing, Mobile Application, Machine Learning, VR/AR Development

\textbf{Languages}:
Python, C, C\#, JavaScript, PHP, Swift, Objective-C, Kotlin, Java, Arduino, Processing, Dart

\textbf{Frameworks \& Libraries}:
TensorFlow, PyTorch, Next.js, Flutter, React, Unity

\textbf{Platforms \& Tools}:
Android, iOS, Web (Frontend/Backend/Database), macOS, Arduino, Raspberry Pi, VR

\textbf{Design \& Prototyping}:
Adobe (Illustrator, Photoshop, Premiere Pro, After Effects, Audition), Figma, Fusion 360
\end{rSection}

%----------------------------------------------------------------------------------------
%   RESEARCH TOPIC SECTION
%----------------------------------------------------------------------------------------
\begin{rSection}{Research Interests}
Human-Computer Interaction (HCI); Interaction Techniques for Mobile Devices; Gesture Design; Sensing Techniques; Machine Learning; 3D Printing; UI/UX Research Methods.
\end{rSection}

%----------------------------------------------------------------------------------------
%   EXPERIENCE & PROJECTS (COMBINED) SECTION
%----------------------------------------------------------------------------------------
\begin{rSection}{Selected Experience \& Projects}

\textbf{LY Corporation --- Research Intern} \hfill \textit{Jul 2024 -- Mar 2025} \\
\textit{(Kotlin, Python, PyTorch, 3D Printing)} \\
During this internship, I carried out advanced research on accessory-based smartphone interfaces, extending the concept of RingSense to new form factors. I investigated IMU and touchscreen-based gesture interactions, implemented CNN models to evaluate recognition accuracy, and built a custom Android app for data collection. Additionally, I fabricated prototypes using a 3D printer and organized user workshops, employing open-coding to identify usage scenarios. This end-to-end experience allowed me to combine hardware prototyping, software development, and user-centered research methods.

\textbf{Personal Finance Web Application} \hfill \textit{Jan 2025 -- Present} \\
\textit{(JavaScript, Firebase, Google Authentication)} \\
I created a web application to manage personal and small business finances, initially tailored to my family's needs. By integrating Firebase for secure data management and Google Authentication for login, I enabled straightforward budget tracking and data visualization. Through this self-directed project, I deepened my understanding of both frontend and backend web technologies. (\href{https://budget-app-db-225db.web.app/}{Link})

\textbf{All Japan Judo Federation --- Freelance Multimedia Editor} \hfill \textit{Dec 2024} \\
\textit{(Adobe Premiere, After Effects)} \\
I produced standardized video-editing templates for the federation's podcast channel, enhancing audience engagement and achieving thousands of YouTube views. This role involved creating efficient workflows with Adobe tools and ensuring consistent branding across multiple video episodes. (\href{https://ltd.japan-judo.org/ajjf-podcast/}{HP}, \href{https://www.youtube.com/playlist?list=PLchwykcjHz9CD9bmAjxs1Xiej3CS-0XqD}{Youtube})

\textbf{HOKUDAI TECH GARAGE: VR Project} \hfill \textit{Mar -- Apr 2022} \\
\textit{(C\#, Unity, Oculus Quest 2)} \\
I developed a VR application that simulates unique animal sensory experiences (e.g., snake pit organs) for the Oculus Quest 2. By conducting user interviews at a local zoo, I refined the system to better address educational needs. This project showcased my ability to deliver a complete user-centered development process, from hardware and software setup to iterative design based on real-world feedback.

\textbf{Gifu City Mayor PR Video} \hfill \textit{2021} \\
\textit{(Adobe Premiere, After Effects)} \\
I produced a promotional video that highlighted local appeal and the Mayor's campaign initiatives. Handling everything from filming and editing to post-production, I iterated the final output through direct client feedback, ensuring the message aligned with the city's goals.

\textbf{Personal ML for Stock Price Prediction} \hfill \textit{2020} \\
\textit{(Python, LightGBM, PCA, ICA)} \\
I built a stock price prediction system using LightGBM and incorporated feature extraction techniques such as principal component analysis (PCA) and independent component analysis (ICA). This allowed me to refine predictions in a real-world context, demonstrating my capacity to integrate multiple ML methods effectively.

\end{rSection}

%----------------------------------------------------------------------------------------
%   AWARDS & FELLOWSHIPS SECTION
%----------------------------------------------------------------------------------------
\begin{rSection}{Fellowships \& Awards}
1. Hokkaido University EXEX Doctoral Fellowship (2024--2027) - \$3,000 annual + stipend \\
2. HOKUDAI TECH GARAGE Project Grant (2022) - \$2,000 per project
\end{rSection}

%----------------------------------------------------------------------------------------
%   EDUCATION SECTION
%----------------------------------------------------------------------------------------
\begin{rSection}{Education}

\textbf{Hokkaido University, Japan} \hfill \textit{April 2024 --- (Expected)}\\
Doctor of Computer Science \quad Advisor: Prof. Daisuke Sakamoto

\textbf{Hokkaido University, Japan} \hfill \textit{April 2022 -- March 2024}\\
Master of Computer Science \quad Advisor: Prof. Daisuke Sakamoto

\textbf{Hokkaido University, Japan} \hfill \textit{April 2020 -- March 2022}\\
Bachelor of Computer Science \quad Advisor: Prof. Daisuke Sakamoto

\textbf{National Institute of Technology (KOSEN), Gifu College, Japan} \hfill \textit{April 2015 -- March 2020}\\
Associate Degree of Engineering \quad Advisor: Prof. Toshinori Deguchi

\end{rSection}

%----------------------------------------------------------------------------------------
%   PUBLICATIONS SECTION
%----------------------------------------------------------------------------------------
\begin{rSection}{Publications}

%--- International Conferences
\begin{rSubsection}{Peer-reviewed International Conferences}{}{}{}
    \item Yuki Abe*, \underline{\textbf{Kan Kusakabe*}}, Myungguen Choi*, Daisuke Sakamoto, Tetsuo Ono. (*Joint first authors).
    \textit{Understanding Usability of VR Pointing Methods with a Handheld-style HMD for Onsite Exhibitions.}
    \textbf{ACM CHI 2025}, April 2025. Honorable Mention Award (\textbf{Acceptance rate 5\%}).
\end{rSubsection}

%--- Domestic Journals
\begin{rSubsection}{Peer-reviewed Domestic Journals}{}{}{}
    \item \underline{\textbf{Kan Kusakabe}}, Daisuke Sakamoto, Tetsuo Ono.
    \textit{RingSense: Exploring User-defined Gestures for Phone Ring Holders}.
    情報処理学会論文誌（IPSJ）, Vol. 66 (2), February 2025.\\
    \textbf{-- Hardware \& ML Details:} Developed a smartphone ring accessory embedding a microcontroller and pressure sensors. Implemented a CNN-based gesture recognition system. Conducted user-defined gesture design (UDG) studies for intuitive interactions.
\end{rSubsection}

%--- Domestic Conferences
\begin{rSubsection}{Peer-reviewed Domestic Conferences}{}{}{}
    \item \underline{\textbf{Kan Kusakabe}}, Daisuke Sakamoto, Tetsuo Ono.
    \textit{スマートフォン背面のジェスチャ入力を実現するスマホリング型デバイスの設計と実装}.
    第30回インタラクティブシステムとソフトウェアに関するワークショップ（WISS 2022）, December 2022. (Long-press, \textbf{Acceptance rate: 12.5\%}).
\end{rSubsection}

%--- Other Publications
\begin{rSubsection}{Other Publications (Non-peer-reviewed)}{}{}{}
    \item \underline{\textbf{Kan Kusakabe*}}, Yuki Abe*, Daisuke Sakamoto, Tetsuo Ono. (*Joint first author).
    \textit{Game-2-X: 種類が異なるゲームプレイ間を繋ぐシステムの提案}.
    第31回インタラクティブシステムとソフトウェアに関するワークショップ（WISS 2023）, デモ発表, December 2023.\\
    \textbf{-- Concept:} Leveraged a reinforcement learning approach to “understand” game states and generate a cross-game scoring mechanism, allowing collaborative or competitive play across different game genres by vectorizing game states.\\

    \item \underline{\textbf{Kan Kusakabe}}, Myungguen Choi, Daisuke Sakamoto, Tetsuo Ono.
    \textit{無段階調整インタフェースのためのハンドジェスチャによる操作手法の探索的研究}. 情報処理学会 研究報告HCI, 2022-HCI-197(25), 1-8, March 2022.\\
\end{rSubsection}

\end{rSection}

\end{document}
